% 03_caches_mmu.tex -- memory hierarchy, inside a 4-way cache, MESI coherence, TLB + page walk.
% Build: tectonic 03_caches_mmu.tex
\documentclass[tikz,border=10pt,12pt]{standalone}
% soc_style.tex -- shared look for the SoC figures (same palette as the slide decks and the
% core register placement figure). Every figure does % soc_style.tex -- shared look for the SoC figures (same palette as the slide decks and the
% core register placement figure). Every figure does % soc_style.tex -- shared look for the SoC figures (same palette as the slide decks and the
% core register placement figure). Every figure does \input{soc_style}.
\usepackage[T1]{fontenc}
\usepackage[scaled=0.95]{helvet}
\renewcommand{\familydefault}{\sfdefault}
\usepackage{booktabs}
\usepackage{amsmath}
\usepackage{amssymb}
\usetikzlibrary{arrows.meta,calc,positioning,fit,shapes.geometric,shapes.misc}

% ---------------- palette ----------------
\definecolor{pcb}{HTML}{0B2E33}
\definecolor{pcbmid}{HTML}{1F5C61}
\definecolor{tint}{HTML}{E6F0EF}
\definecolor{copper}{HTML}{B5651D}
\definecolor{coppert}{HTML}{F6E8DC}
\definecolor{ink}{HTML}{1F2A2E}
\definecolor{muted}{HTML}{5B6B70}
\definecolor{hardf}{HTML}{DCE6F2}
\definecolor{harde}{HTML}{2F4A7A}
\definecolor{good}{HTML}{2E7D32}
\definecolor{bad}{HTML}{B3261E}
% one colour per kind of block: border (d) / fill (f)
\definecolor{cpud}{HTML}{2B5C9E} \definecolor{cpuf}{HTML}{D6E4F5}
\definecolor{gpud}{HTML}{2E7D32} \definecolor{gpuf}{HTML}{D5EDD0}
\definecolor{npud}{HTML}{0B2E33} \definecolor{npuf}{HTML}{BFE0DB}
\definecolor{memd}{HTML}{5E3A8C} \definecolor{memf}{HTML}{E4DAF3}
\definecolor{busd}{HTML}{B5651D} \definecolor{busf}{HTML}{F6E8DC}
\definecolor{prd}{HTML}{546E7A}  \definecolor{prf}{HTML}{E3EAED}
\definecolor{secd}{HTML}{B3261E} \definecolor{secf}{HTML}{F7D9D6}
\definecolor{pwrd}{HTML}{8A6800} \definecolor{pwrf}{HTML}{FFF0B8}
\definecolor{medd}{HTML}{8E2463} \definecolor{medf}{HTML}{F5D6E6}
\definecolor{skd}{HTML}{2E7D32}  \definecolor{skf}{HTML}{BFE3B4}   % skew buffers, as in the register figure

\tikzset{
  blk/.style 2 args={draw=#1, fill=#2, rounded corners=4pt, thick, align=center,
                     font=\scriptsize, inner sep=3pt},
  cpu/.style={blk={cpud}{cpuf}},  gpu/.style={blk={gpud}{gpuf}},
  npu/.style={blk={npud}{npuf}},  mem/.style={blk={memd}{memf}},
  bus/.style={blk={busd}{busf}},  per/.style={blk={prd}{prf}},
  sec/.style={blk={secd}{secf}},  pwr/.style={blk={pwrd}{pwrf}},
  med/.style={blk={medd}{medf}},
  comb/.style ={draw=copper, fill=coppert, rounded corners=2pt, thick, align=center,
                font=\scriptsize, inner sep=2pt},                       % combinational logic
  hard/.style ={draw=harde, fill=hardf, thick, dashed, align=center,
                font=\scriptsize, inner sep=2pt},                       % hard macro / off-chip
  sram/.style ={draw=memd, fill=memf, thick, align=center, font=\scriptsize,
                inner sep=2pt},                                         % SRAM array, not flip-flops
  ff/.style   ={draw=#1, fill=#1!30, line width=0.35pt},                % one flip-flop square
  ffgone/.style={draw=muted, fill=white, densely dashed, line width=0.35pt},
  panel/.style={draw=muted, rounded corners=6pt, line width=0.9pt},
  ptitle/.style={anchor=north west, font=\normalsize\bfseries, text=ink, inner sep=0pt},
  note/.style ={font=\scriptsize, text=ink, align=left, inner sep=0pt},
  lab/.style  ={font=\tiny, text=muted, inner sep=1pt},
  axi/.style  ={{Stealth[length=2.6mm]}-{Stealth[length=2.6mm]}, line width=2.2pt, copper},
  axi1/.style ={-{Stealth[length=2.6mm]}, line width=2.2pt, copper},
  apb/.style  ={{Stealth[length=1.8mm]}-{Stealth[length=1.8mm]}, line width=1pt, copper},
  apb1/.style ={-{Stealth[length=1.8mm]}, line width=1pt, copper},
  dat/.style  ={-{Stealth[length=1.8mm]}, thick, pcbmid},
  ctl/.style  ={-{Stealth[length=1.8mm]}, densely dashed, thick, muted},
  irq/.style  ={-{Stealth[length=2mm]}, densely dotted, line width=1.3pt, bad},
  num/.style  ={circle, fill=copper, text=white, font=\scriptsize\bfseries, inner sep=1pt,
                minimum size=0.46cm},
  ours/.style ={draw=bad, line width=1.2pt, densely dashed, rounded corners=4pt},
}

% ---------------- helpers ----------------
\newcommand{\rn}[1]{\texttt{#1}}
% \ffrow{x}{y}{n}{cell}{style}: n flip-flops in a row, bottom-left corner at (x,y)
\newcommand{\ffrow}[5]{\foreach \q in {1,...,#3}{\draw[#5] ({#1+(\q-1)*#4},{#2}) rectangle ++(#4,#4);}}
% \ffcol{x}{y}{n}{cell}{style}: n flip-flops in a column, top-left corner at (x,y)
\newcommand{\ffcol}[5]{\foreach \q in {1,...,#3}{\draw[#5] ({#1},{#2-\q*#4}) rectangle ++(#4,#4);}}
% \ffgrid{x}{ytop}{n}{per row}{cell}{style}: n flip-flops, filled row by row from the top-left
\newcommand{\ffgrid}[6]{\pgfmathtruncatemacro{\gl}{#3-1}%
  \foreach \b in {0,...,\gl}{\pgfmathtruncatemacro{\gc}{mod(\b,#4)}\pgfmathtruncatemacro{\gr}{div(\b,#4)}%
    \draw[#6] ({#1+\gc*#5},{#2-(\gr+1)*#5}) rectangle ++(#5,#5);}}

% every figure: same page size, title, footer
\newcommand{\canvas}{\useasboundingbox (-0.3,-0.2) rectangle (38.6,30.4);}
\newcommand{\figtitle}[3]{%
  \node[num, minimum size=0.9cm, font=\Large\bfseries] at (0.45,29.55) {#1};
  \node[anchor=west, font=\LARGE\bfseries, text=pcb] at (1.15,29.55) {#2};
  \node[anchor=west, font=\small, text=ink] at (0,28.7) {#3};}
\newcommand{\figfoot}[1]{\node[anchor=west, font=\scriptsize, text=muted] at (0,0.25) {#1};}
% \panel{x0}{y0}{x1}{y1}{title}
\newcommand{\panel}[5]{\draw[panel] (#1,#2) rectangle (#3,#4);
  \node[ptitle] at ({#1+0.3},{#4-0.25}) {#5};}

% ---------------- timing diagrams: 1 cycle = \cw cm, low at y, high at y+0.45 ----------------
\newcommand{\cw}{1.0}
\newcommand{\wclk}[3]{% x0 y ncycles
  \foreach \k in {1,...,#3}{\draw[pcb, thick] ({#1+(\k-1)*\cw},#2) -- ++(0,0.45) -- ++(0.5*\cw,0)
     -- ++(0,-0.45) -- ++(0.5*\cw,0);}}
\newcommand{\wsig}[3]{% x0 y {level per cycle, e.g. 0,1,1,0}
  \gdef\wprev{x}%
  \foreach \v [count=\k from 0] in {#3}{%
    \draw[pcb, thick] ({#1+\k*\cw},{#2+0.45*\v}) -- ++(\cw,0);
    \if x\wprev\else\ifnum\wprev=\v\else\draw[pcb, thick] ({#1+\k*\cw},#2) -- ++(0,0.45);\fi\fi
    \global\let\wprev\v}}
\newcommand{\wbus}[5]{% x0 y c0 c1 text: one value held from cycle c0 to cycle c1
  \draw[pcbmid, thick, fill=white] ({#1+(#3)*\cw},{#2+0.225}) -- ({#1+(#3)*\cw+0.08},{#2+0.45})
    -- ({#1+(#4)*\cw-0.08},{#2+0.45}) -- ({#1+(#4)*\cw},{#2+0.225}) -- ({#1+(#4)*\cw-0.08},#2)
    -- ({#1+(#3)*\cw+0.08},#2) -- cycle;
  \node[font=\tiny\ttfamily] at ({#1+((#3)+(#4))*\cw/2},{#2+0.225}) {#5};}
\newcommand{\widle}[4]{% x0 y c0 c1: bus idle (no value)
  \draw[muted, thick] ({#1+(#3)*\cw},{#2+0.225}) -- ({#1+(#4)*\cw},{#2+0.225});}
\newcommand{\wname}[3]{\node[anchor=east, font=\scriptsize\ttfamily] at ({#1-0.15},{#2+0.225}) {#3};}
.
\usepackage[T1]{fontenc}
\usepackage[scaled=0.95]{helvet}
\renewcommand{\familydefault}{\sfdefault}
\usepackage{booktabs}
\usepackage{amsmath}
\usepackage{amssymb}
\usetikzlibrary{arrows.meta,calc,positioning,fit,shapes.geometric,shapes.misc}

% ---------------- palette ----------------
\definecolor{pcb}{HTML}{0B2E33}
\definecolor{pcbmid}{HTML}{1F5C61}
\definecolor{tint}{HTML}{E6F0EF}
\definecolor{copper}{HTML}{B5651D}
\definecolor{coppert}{HTML}{F6E8DC}
\definecolor{ink}{HTML}{1F2A2E}
\definecolor{muted}{HTML}{5B6B70}
\definecolor{hardf}{HTML}{DCE6F2}
\definecolor{harde}{HTML}{2F4A7A}
\definecolor{good}{HTML}{2E7D32}
\definecolor{bad}{HTML}{B3261E}
% one colour per kind of block: border (d) / fill (f)
\definecolor{cpud}{HTML}{2B5C9E} \definecolor{cpuf}{HTML}{D6E4F5}
\definecolor{gpud}{HTML}{2E7D32} \definecolor{gpuf}{HTML}{D5EDD0}
\definecolor{npud}{HTML}{0B2E33} \definecolor{npuf}{HTML}{BFE0DB}
\definecolor{memd}{HTML}{5E3A8C} \definecolor{memf}{HTML}{E4DAF3}
\definecolor{busd}{HTML}{B5651D} \definecolor{busf}{HTML}{F6E8DC}
\definecolor{prd}{HTML}{546E7A}  \definecolor{prf}{HTML}{E3EAED}
\definecolor{secd}{HTML}{B3261E} \definecolor{secf}{HTML}{F7D9D6}
\definecolor{pwrd}{HTML}{8A6800} \definecolor{pwrf}{HTML}{FFF0B8}
\definecolor{medd}{HTML}{8E2463} \definecolor{medf}{HTML}{F5D6E6}
\definecolor{skd}{HTML}{2E7D32}  \definecolor{skf}{HTML}{BFE3B4}   % skew buffers, as in the register figure

\tikzset{
  blk/.style 2 args={draw=#1, fill=#2, rounded corners=4pt, thick, align=center,
                     font=\scriptsize, inner sep=3pt},
  cpu/.style={blk={cpud}{cpuf}},  gpu/.style={blk={gpud}{gpuf}},
  npu/.style={blk={npud}{npuf}},  mem/.style={blk={memd}{memf}},
  bus/.style={blk={busd}{busf}},  per/.style={blk={prd}{prf}},
  sec/.style={blk={secd}{secf}},  pwr/.style={blk={pwrd}{pwrf}},
  med/.style={blk={medd}{medf}},
  comb/.style ={draw=copper, fill=coppert, rounded corners=2pt, thick, align=center,
                font=\scriptsize, inner sep=2pt},                       % combinational logic
  hard/.style ={draw=harde, fill=hardf, thick, dashed, align=center,
                font=\scriptsize, inner sep=2pt},                       % hard macro / off-chip
  sram/.style ={draw=memd, fill=memf, thick, align=center, font=\scriptsize,
                inner sep=2pt},                                         % SRAM array, not flip-flops
  ff/.style   ={draw=#1, fill=#1!30, line width=0.35pt},                % one flip-flop square
  ffgone/.style={draw=muted, fill=white, densely dashed, line width=0.35pt},
  panel/.style={draw=muted, rounded corners=6pt, line width=0.9pt},
  ptitle/.style={anchor=north west, font=\normalsize\bfseries, text=ink, inner sep=0pt},
  note/.style ={font=\scriptsize, text=ink, align=left, inner sep=0pt},
  lab/.style  ={font=\tiny, text=muted, inner sep=1pt},
  axi/.style  ={{Stealth[length=2.6mm]}-{Stealth[length=2.6mm]}, line width=2.2pt, copper},
  axi1/.style ={-{Stealth[length=2.6mm]}, line width=2.2pt, copper},
  apb/.style  ={{Stealth[length=1.8mm]}-{Stealth[length=1.8mm]}, line width=1pt, copper},
  apb1/.style ={-{Stealth[length=1.8mm]}, line width=1pt, copper},
  dat/.style  ={-{Stealth[length=1.8mm]}, thick, pcbmid},
  ctl/.style  ={-{Stealth[length=1.8mm]}, densely dashed, thick, muted},
  irq/.style  ={-{Stealth[length=2mm]}, densely dotted, line width=1.3pt, bad},
  num/.style  ={circle, fill=copper, text=white, font=\scriptsize\bfseries, inner sep=1pt,
                minimum size=0.46cm},
  ours/.style ={draw=bad, line width=1.2pt, densely dashed, rounded corners=4pt},
}

% ---------------- helpers ----------------
\newcommand{\rn}[1]{\texttt{#1}}
% \ffrow{x}{y}{n}{cell}{style}: n flip-flops in a row, bottom-left corner at (x,y)
\newcommand{\ffrow}[5]{\foreach \q in {1,...,#3}{\draw[#5] ({#1+(\q-1)*#4},{#2}) rectangle ++(#4,#4);}}
% \ffcol{x}{y}{n}{cell}{style}: n flip-flops in a column, top-left corner at (x,y)
\newcommand{\ffcol}[5]{\foreach \q in {1,...,#3}{\draw[#5] ({#1},{#2-\q*#4}) rectangle ++(#4,#4);}}
% \ffgrid{x}{ytop}{n}{per row}{cell}{style}: n flip-flops, filled row by row from the top-left
\newcommand{\ffgrid}[6]{\pgfmathtruncatemacro{\gl}{#3-1}%
  \foreach \b in {0,...,\gl}{\pgfmathtruncatemacro{\gc}{mod(\b,#4)}\pgfmathtruncatemacro{\gr}{div(\b,#4)}%
    \draw[#6] ({#1+\gc*#5},{#2-(\gr+1)*#5}) rectangle ++(#5,#5);}}

% every figure: same page size, title, footer
\newcommand{\canvas}{\useasboundingbox (-0.3,-0.2) rectangle (38.6,30.4);}
\newcommand{\figtitle}[3]{%
  \node[num, minimum size=0.9cm, font=\Large\bfseries] at (0.45,29.55) {#1};
  \node[anchor=west, font=\LARGE\bfseries, text=pcb] at (1.15,29.55) {#2};
  \node[anchor=west, font=\small, text=ink] at (0,28.7) {#3};}
\newcommand{\figfoot}[1]{\node[anchor=west, font=\scriptsize, text=muted] at (0,0.25) {#1};}
% \panel{x0}{y0}{x1}{y1}{title}
\newcommand{\panel}[5]{\draw[panel] (#1,#2) rectangle (#3,#4);
  \node[ptitle] at ({#1+0.3},{#4-0.25}) {#5};}

% ---------------- timing diagrams: 1 cycle = \cw cm, low at y, high at y+0.45 ----------------
\newcommand{\cw}{1.0}
\newcommand{\wclk}[3]{% x0 y ncycles
  \foreach \k in {1,...,#3}{\draw[pcb, thick] ({#1+(\k-1)*\cw},#2) -- ++(0,0.45) -- ++(0.5*\cw,0)
     -- ++(0,-0.45) -- ++(0.5*\cw,0);}}
\newcommand{\wsig}[3]{% x0 y {level per cycle, e.g. 0,1,1,0}
  \gdef\wprev{x}%
  \foreach \v [count=\k from 0] in {#3}{%
    \draw[pcb, thick] ({#1+\k*\cw},{#2+0.45*\v}) -- ++(\cw,0);
    \if x\wprev\else\ifnum\wprev=\v\else\draw[pcb, thick] ({#1+\k*\cw},#2) -- ++(0,0.45);\fi\fi
    \global\let\wprev\v}}
\newcommand{\wbus}[5]{% x0 y c0 c1 text: one value held from cycle c0 to cycle c1
  \draw[pcbmid, thick, fill=white] ({#1+(#3)*\cw},{#2+0.225}) -- ({#1+(#3)*\cw+0.08},{#2+0.45})
    -- ({#1+(#4)*\cw-0.08},{#2+0.45}) -- ({#1+(#4)*\cw},{#2+0.225}) -- ({#1+(#4)*\cw-0.08},#2)
    -- ({#1+(#3)*\cw+0.08},#2) -- cycle;
  \node[font=\tiny\ttfamily] at ({#1+((#3)+(#4))*\cw/2},{#2+0.225}) {#5};}
\newcommand{\widle}[4]{% x0 y c0 c1: bus idle (no value)
  \draw[muted, thick] ({#1+(#3)*\cw},{#2+0.225}) -- ({#1+(#4)*\cw},{#2+0.225});}
\newcommand{\wname}[3]{\node[anchor=east, font=\scriptsize\ttfamily] at ({#1-0.15},{#2+0.225}) {#3};}
.
\usepackage[T1]{fontenc}
\usepackage[scaled=0.95]{helvet}
\renewcommand{\familydefault}{\sfdefault}
\usepackage{booktabs}
\usepackage{amsmath}
\usepackage{amssymb}
\usetikzlibrary{arrows.meta,calc,positioning,fit,shapes.geometric,shapes.misc}

% ---------------- palette ----------------
\definecolor{pcb}{HTML}{0B2E33}
\definecolor{pcbmid}{HTML}{1F5C61}
\definecolor{tint}{HTML}{E6F0EF}
\definecolor{copper}{HTML}{B5651D}
\definecolor{coppert}{HTML}{F6E8DC}
\definecolor{ink}{HTML}{1F2A2E}
\definecolor{muted}{HTML}{5B6B70}
\definecolor{hardf}{HTML}{DCE6F2}
\definecolor{harde}{HTML}{2F4A7A}
\definecolor{good}{HTML}{2E7D32}
\definecolor{bad}{HTML}{B3261E}
% one colour per kind of block: border (d) / fill (f)
\definecolor{cpud}{HTML}{2B5C9E} \definecolor{cpuf}{HTML}{D6E4F5}
\definecolor{gpud}{HTML}{2E7D32} \definecolor{gpuf}{HTML}{D5EDD0}
\definecolor{npud}{HTML}{0B2E33} \definecolor{npuf}{HTML}{BFE0DB}
\definecolor{memd}{HTML}{5E3A8C} \definecolor{memf}{HTML}{E4DAF3}
\definecolor{busd}{HTML}{B5651D} \definecolor{busf}{HTML}{F6E8DC}
\definecolor{prd}{HTML}{546E7A}  \definecolor{prf}{HTML}{E3EAED}
\definecolor{secd}{HTML}{B3261E} \definecolor{secf}{HTML}{F7D9D6}
\definecolor{pwrd}{HTML}{8A6800} \definecolor{pwrf}{HTML}{FFF0B8}
\definecolor{medd}{HTML}{8E2463} \definecolor{medf}{HTML}{F5D6E6}
\definecolor{skd}{HTML}{2E7D32}  \definecolor{skf}{HTML}{BFE3B4}   % skew buffers, as in the register figure

\tikzset{
  blk/.style 2 args={draw=#1, fill=#2, rounded corners=4pt, thick, align=center,
                     font=\scriptsize, inner sep=3pt},
  cpu/.style={blk={cpud}{cpuf}},  gpu/.style={blk={gpud}{gpuf}},
  npu/.style={blk={npud}{npuf}},  mem/.style={blk={memd}{memf}},
  bus/.style={blk={busd}{busf}},  per/.style={blk={prd}{prf}},
  sec/.style={blk={secd}{secf}},  pwr/.style={blk={pwrd}{pwrf}},
  med/.style={blk={medd}{medf}},
  comb/.style ={draw=copper, fill=coppert, rounded corners=2pt, thick, align=center,
                font=\scriptsize, inner sep=2pt},                       % combinational logic
  hard/.style ={draw=harde, fill=hardf, thick, dashed, align=center,
                font=\scriptsize, inner sep=2pt},                       % hard macro / off-chip
  sram/.style ={draw=memd, fill=memf, thick, align=center, font=\scriptsize,
                inner sep=2pt},                                         % SRAM array, not flip-flops
  ff/.style   ={draw=#1, fill=#1!30, line width=0.35pt},                % one flip-flop square
  ffgone/.style={draw=muted, fill=white, densely dashed, line width=0.35pt},
  panel/.style={draw=muted, rounded corners=6pt, line width=0.9pt},
  ptitle/.style={anchor=north west, font=\normalsize\bfseries, text=ink, inner sep=0pt},
  note/.style ={font=\scriptsize, text=ink, align=left, inner sep=0pt},
  lab/.style  ={font=\tiny, text=muted, inner sep=1pt},
  axi/.style  ={{Stealth[length=2.6mm]}-{Stealth[length=2.6mm]}, line width=2.2pt, copper},
  axi1/.style ={-{Stealth[length=2.6mm]}, line width=2.2pt, copper},
  apb/.style  ={{Stealth[length=1.8mm]}-{Stealth[length=1.8mm]}, line width=1pt, copper},
  apb1/.style ={-{Stealth[length=1.8mm]}, line width=1pt, copper},
  dat/.style  ={-{Stealth[length=1.8mm]}, thick, pcbmid},
  ctl/.style  ={-{Stealth[length=1.8mm]}, densely dashed, thick, muted},
  irq/.style  ={-{Stealth[length=2mm]}, densely dotted, line width=1.3pt, bad},
  num/.style  ={circle, fill=copper, text=white, font=\scriptsize\bfseries, inner sep=1pt,
                minimum size=0.46cm},
  ours/.style ={draw=bad, line width=1.2pt, densely dashed, rounded corners=4pt},
}

% ---------------- helpers ----------------
\newcommand{\rn}[1]{\texttt{#1}}
% \ffrow{x}{y}{n}{cell}{style}: n flip-flops in a row, bottom-left corner at (x,y)
\newcommand{\ffrow}[5]{\foreach \q in {1,...,#3}{\draw[#5] ({#1+(\q-1)*#4},{#2}) rectangle ++(#4,#4);}}
% \ffcol{x}{y}{n}{cell}{style}: n flip-flops in a column, top-left corner at (x,y)
\newcommand{\ffcol}[5]{\foreach \q in {1,...,#3}{\draw[#5] ({#1},{#2-\q*#4}) rectangle ++(#4,#4);}}
% \ffgrid{x}{ytop}{n}{per row}{cell}{style}: n flip-flops, filled row by row from the top-left
\newcommand{\ffgrid}[6]{\pgfmathtruncatemacro{\gl}{#3-1}%
  \foreach \b in {0,...,\gl}{\pgfmathtruncatemacro{\gc}{mod(\b,#4)}\pgfmathtruncatemacro{\gr}{div(\b,#4)}%
    \draw[#6] ({#1+\gc*#5},{#2-(\gr+1)*#5}) rectangle ++(#5,#5);}}

% every figure: same page size, title, footer
\newcommand{\canvas}{\useasboundingbox (-0.3,-0.2) rectangle (38.6,30.4);}
\newcommand{\figtitle}[3]{%
  \node[num, minimum size=0.9cm, font=\Large\bfseries] at (0.45,29.55) {#1};
  \node[anchor=west, font=\LARGE\bfseries, text=pcb] at (1.15,29.55) {#2};
  \node[anchor=west, font=\small, text=ink] at (0,28.7) {#3};}
\newcommand{\figfoot}[1]{\node[anchor=west, font=\scriptsize, text=muted] at (0,0.25) {#1};}
% \panel{x0}{y0}{x1}{y1}{title}
\newcommand{\panel}[5]{\draw[panel] (#1,#2) rectangle (#3,#4);
  \node[ptitle] at ({#1+0.3},{#4-0.25}) {#5};}

% ---------------- timing diagrams: 1 cycle = \cw cm, low at y, high at y+0.45 ----------------
\newcommand{\cw}{1.0}
\newcommand{\wclk}[3]{% x0 y ncycles
  \foreach \k in {1,...,#3}{\draw[pcb, thick] ({#1+(\k-1)*\cw},#2) -- ++(0,0.45) -- ++(0.5*\cw,0)
     -- ++(0,-0.45) -- ++(0.5*\cw,0);}}
\newcommand{\wsig}[3]{% x0 y {level per cycle, e.g. 0,1,1,0}
  \gdef\wprev{x}%
  \foreach \v [count=\k from 0] in {#3}{%
    \draw[pcb, thick] ({#1+\k*\cw},{#2+0.45*\v}) -- ++(\cw,0);
    \if x\wprev\else\ifnum\wprev=\v\else\draw[pcb, thick] ({#1+\k*\cw},#2) -- ++(0,0.45);\fi\fi
    \global\let\wprev\v}}
\newcommand{\wbus}[5]{% x0 y c0 c1 text: one value held from cycle c0 to cycle c1
  \draw[pcbmid, thick, fill=white] ({#1+(#3)*\cw},{#2+0.225}) -- ({#1+(#3)*\cw+0.08},{#2+0.45})
    -- ({#1+(#4)*\cw-0.08},{#2+0.45}) -- ({#1+(#4)*\cw},{#2+0.225}) -- ({#1+(#4)*\cw-0.08},#2)
    -- ({#1+(#3)*\cw+0.08},#2) -- cycle;
  \node[font=\tiny\ttfamily] at ({#1+((#3)+(#4))*\cw/2},{#2+0.225}) {#5};}
\newcommand{\widle}[4]{% x0 y c0 c1: bus idle (no value)
  \draw[muted, thick] ({#1+(#3)*\cw},{#2+0.225}) -- ({#1+(#4)*\cw},{#2+0.225});}
\newcommand{\wname}[3]{\node[anchor=east, font=\scriptsize\ttfamily] at ({#1-0.15},{#2+0.225}) {#3};}

\begin{document}
\begin{tikzpicture}[line join=round]
\canvas
\figtitle{3}{Caches and MMU: how the CPU finds its data fast}
  {A cache is a small SRAM that keeps copies of recently used memory. The address is cut into tag, index and offset;
   comparators decide hit or miss in one cycle. One example address is followed through the whole page.}

% =====================================================================
% (A) memory hierarchy
% =====================================================================
\panel{0}{13.6}{10.5}{28.1}{Memory hierarchy (typical)}
\foreach \nm/\col [count=\k from 0] in {registers/cpuf, L1 cache/cpuf, L2 cache/cpuf, {L3 / SLC}/memf, DRAM/memf, flash/hardf}{
  \pgfmathsetmacro{\yt}{26.6-\k*1.2}\pgfmathsetmacro{\yb}{\yt-1.2}
  \pgfmathsetmacro{\wt}{1.0+1.5*\k}\pgfmathsetmacro{\wb}{2.5+1.5*\k}
  \draw[pcb, thick, fill=\col] ({5.25-\wt/2},\yt) -- ({5.25+\wt/2},\yt) -- ({5.25+\wb/2},\yb) -- ({5.25-\wb/2},\yb) -- cycle;
  \node[font=\scriptsize\bfseries] at (5.25,{(\yt+\yb)/2}) {\nm};}
\draw[-{Stealth[length=2mm]}, thick, good] (0.45,19.6) -- (0.45,26.4);
\node[font=\tiny\bfseries, text=good, rotate=90, anchor=south] at (0.4,23.0) {faster, smaller, costlier per bit};
\node[anchor=north west, font=\scriptsize, inner sep=0pt] at (0.35,19.0) {%
  \renewcommand{\arraystretch}{1.1}%
  \begin{tabular}{@{}l l l@{}}
    \toprule
    \textbf{Level} & \textbf{Size} & \textbf{Time to read}\\
    \midrule
    registers  & 32 $\times$ 32 bit (RV32) & inside the cycle\\
    L1 (per core) & 32--64 KB          & $\approx$4 cycles\\
    L2         & 256 KB -- 2 MB        & $\approx$10--15 cycles\\
    L3 / SLC   & 4--32 MB, shared      & $\approx$30--50 cycles\\
    DRAM       & 4--16 GB              & $\approx$100 ns $\approx$ 300 cycles\\
    flash      & 128 GB -- 1 TB        & $\approx$50--100 $\mu$s\\
    \bottomrule
  \end{tabular}};
\node[note, font=\small, anchor=north west, text width=9.8cm] at (0.35,15.7) {%
  Works because programs reuse recent data and use nearby data. Most loads hit in L1.};

% =====================================================================
% (B) inside one L1 data cache
% =====================================================================
\panel{10.8}{13.6}{38.3}{28.1}{Inside one L1 data cache: 32 KB, 4 ways, 64-byte lines}
% address bar
\node[anchor=west, font=\scriptsize\bfseries] at (11.1,26.85) {physical address 0x8000\_1A44 =};
\foreach \b in {0,...,31}{
  \pgfmathsetmacro{\x}{16.0+(31-\b)*0.36}
  \ifnum\b>12 \def\fc{cpuf}\else\ifnum\b>5 \def\fc{pwrf}\else\def\fc{medf}\fi\fi
  \draw[pcb, fill=\fc] (\x,26.65) rectangle ++(0.36,0.4);}
\draw[pcb, very thick] (16.0,26.65) rectangle (27.52,27.05);
\node[font=\scriptsize\bfseries, text=cpud] at (19.42,26.4) {tag [31:13] = 19 bit = 0x40000};
\node[font=\scriptsize\bfseries, text=pwrd] at (24.82,26.4) {index [12:6] = 105};
\node[font=\scriptsize\bfseries, text=medd] at (26.44,27.25) {offset [5:0] = 4};
\node[note, anchor=west, text width=10.2cm] at (27.9,26.85)
  {index picks 1 of 128 sets $\cdot$ tag says which memory line it is $\cdot$ offset picks the byte in the 64-byte line};

% index decoder
\node[comb, minimum width=1.25cm, minimum height=3.6cm, text width=1.1cm, anchor=south west, font=\tiny] at (11.05,21.9)
  {\textbf{index decoder}\\7 bit $\rightarrow$ 1 of 128 rows};
\draw[dat] (24.82,26.15) -- (24.82,25.95) -- (11.72,25.95) -- (11.72,25.5);

% four ways
\foreach \w/\v/\d/\t/\res in {0/1/0/0x1F3A2/{$\neq$}, 1/0/0/0x40000/{V=0}, 2/1/1/0x40000/{\textbf{=}}, 3/1/0/0x2C410/{$\neq$}}{
  \pgfmathsetmacro{\x}{13.5+\w*5.0}
  \node[font=\scriptsize\bfseries, anchor=south west] at (\x,25.35) {way \w};
  \node[lab, anchor=south] at (\x+0.15,25.0) {V};
  \node[lab, anchor=south] at (\x+0.45,25.0) {D};
  \node[lab, anchor=south] at (\x+1.35,25.0) {tag (19)};
  \node[lab, anchor=south] at (\x+3.46,25.0) {data: 64 B = 16 words};
  \foreach \row/\yy in {set 0/24.6, set 1/24.2, {\vdots}/23.8, set 105/23.4, {\vdots}/23.0, set 127/22.6}{
    \draw[pcb, fill=white] (\x,\yy-0.17) rectangle ++(0.3,0.36);
    \draw[pcb, fill=white] (\x+0.3,\yy-0.17) rectangle ++(0.3,0.36);
    \draw[cpud, fill=cpuf] (\x+0.6,\yy-0.17) rectangle ++(1.5,0.36);
    \foreach \q in {0,...,15}{\draw[memd, fill=memf] ({\x+2.1+\q*0.17},\yy-0.17) rectangle ++(0.17,0.36);}}
  \draw[bad, line width=1.4pt] (\x-0.05,23.18) rectangle (\x+4.87,23.64);
  \node[font=\tiny\bfseries] at (\x+0.15,23.41) {\v};
  \node[font=\tiny\bfseries] at (\x+0.45,23.41) {\d};
  \node[font=\tiny\ttfamily\bfseries] at (\x+1.35,23.41) {\t};
  \node[circle, comb, minimum size=0.6cm, inner sep=0pt, font=\scriptsize] (cmp\w) at (\x+1.35,21.45) {\res};
  \draw[dat] (\x+1.35,22.4) -- (cmp\w.north);
  \draw[dat] (\x+0.15,22.4) -- (\x+0.15,21.45) -- (cmp\w.west);
}
\foreach \row/\yy in {set 0/24.6, set 1/24.2, set 105/23.4, set 127/22.6}{\node[lab, anchor=east] at (13.45,\yy) {\row};}
% address tag bus to the comparators
\draw[dat] (16.4,26.65) -- (16.4,26.25) -- (10.95,26.25) -- (10.95,20.75) -- (33.3,20.75);
\foreach \w in {0,...,3}{\pgfmathsetmacro{\x}{13.5+\w*5.0+1.35}\draw[dat] (\x,20.75) -- (\x,21.15);}
\node[lab, anchor=north west, fill=white] at (11.0,20.72) {address tag 0x40000 goes to all 4 comparators at the same time};
\node[note, anchor=north west, text width=4.2cm] at (34.4,26.0) {%
  \textbf{size check}\\
  32 KB / 64 B = 512 lines\\
  512 / 4 ways = 128 sets\\[3pt]
  tag store: 4 $\times$ 128 $\times$ (19+2)\\= 10,752 bits\\
  data store: 262,144 bits\\[3pt]
  both are \textbf{SRAM}, not FF\\
  + 3 bits per set for\\``least recently used''};

% hit logic
\node[comb, minimum width=4.0cm, minimum height=0.85cm, anchor=south west, font=\tiny] (hit) at (12.6,19.1)
  {\textbf{hit} = OR of (tag equal AND V) $\rightarrow$ way 2};
\node[comb, minimum width=5.6cm, minimum height=0.85cm, anchor=south west, font=\tiny] (wm) at (17.3,19.1)
  {\textbf{way mux} 4 $\rightarrow$ 1: pick the 64-byte line of way 2};
\node[comb, minimum width=6.2cm, minimum height=0.85cm, anchor=south west, font=\tiny] (om) at (23.6,19.1)
  {\textbf{word mux} 16 $\rightarrow$ 1: offset[5:2] = 1 $\rightarrow$ bytes 4--7};
\node[cpu, minimum width=3.3cm, minimum height=0.85cm, anchor=south west, font=\tiny] (out) at (30.5,19.1)
  {32-bit data $\rightarrow$ CPU};
\draw[dat] (hit) -- (wm); \draw[dat] (wm) -- (om); \draw[dat] (om) -- (out);

\node[note, anchor=north west, text width=12.6cm] at (11.1,18.45) {%
  \textbf{Hit} (this example): set 105 is read in all 4 ways at once, way 2 has V = 1 and tag 0x40000 $\Rightarrow$ data in $\approx$4 cycles.\\[3pt]
  \textbf{Miss}: ask L2 for the 64-byte line; choose a victim way (least recently used); if its D = 1, write it back first;
  fill the line, then answer the CPU. A miss buffer lets the core continue with other loads meanwhile.};
\node[note, anchor=north west, text width=12.4cm] at (25.3,18.45) {%
  \textbf{Write-back}: a store changes only the cache line and sets D = 1. Memory is updated when the line is thrown out.\\[3pt]
  \textbf{V} = line holds real data. \textbf{D} = line is newer than memory.
  Big caches add 2 state bits per line for coherence (next panel).\\[3pt]
  A hit costs 1 SRAM read + 4 compares of 19 bits + 2 muxes: all inside one pipeline stage group.};

% =====================================================================
% (C) coherence: MESI
% =====================================================================
\panel{0}{0.7}{14.5}{13.3}{Coherence: MESI states of one line}
\node[circle, cpu, minimum size=1.5cm, font=\normalsize\bfseries] (M) at (2.6,10.0) {M};
\node[circle, cpu, minimum size=1.5cm, font=\normalsize\bfseries] (E) at (9.0,10.0) {E};
\node[circle, cpu, minimum size=1.5cm, font=\normalsize\bfseries] (S) at (9.0,4.6) {S};
\node[circle, draw=muted, fill=white, thick, minimum size=1.5cm, font=\normalsize\bfseries] (I) at (2.6,4.6) {I};
\draw[dat] (I) to[bend left=12] node[lab, text=ink, fill=white, pos=0.5] {read, no other copy} (E);
\draw[dat] (I) to[bend right=12] node[lab, text=ink, fill=white, below] {read, others have it} (S);
\draw[dat] (I) to[bend left=20] node[lab, text=ink, fill=white, left] {write} (M);
\draw[dat] (E) to node[lab, text=ink, fill=white, above] {my write (silent)} (M);
\draw[dat] (S) to[bend right=8] node[lab, text=ink, fill=white, pos=0.35] {my write: invalidate others} (M);
\draw[ctl] (M) to[bend right=30] node[lab, text=ink, fill=white, pos=0.6] {other core reads} (S);
\draw[ctl] (E) to node[lab, text=ink, fill=white, right] {other core reads} (S);
\draw[ctl] (S) to[bend left=25] node[lab, text=ink, fill=white, below] {other core writes} (I);
\node[note, anchor=north west, text width=13.8cm] at (0.35,2.85) {%
  \textbf{M} only copy, changed $\cdot$ \textbf{E} only copy, clean $\cdot$ \textbf{S} clean, others may have it $\cdot$ \textbf{I} not valid.
  Solid = this core's action, dashed = seen from another core (snoop).\\[3pt]
  Stored as 2 extra bits next to each tag. A snoop filter in the shared L3 remembers which core has which line, so not every core is asked.};

% =====================================================================
% (D) MMU: TLB + page walk
% =====================================================================
\panel{14.8}{0.7}{38.3}{13.3}{MMU: virtual address $\rightarrow$ physical address (RISC-V Sv32, 4 KB pages)}
\node[anchor=west, font=\scriptsize\bfseries] at (15.1,11.75) {virtual 0x0040\_3A44 =};
\node[cpu, minimum width=3.0cm, minimum height=0.5cm, anchor=west, inner sep=0pt] (v1) at (18.6,11.75) {VPN1 = 0x001 (10)};
\node[cpu, minimum width=3.0cm, minimum height=0.5cm, anchor=west, inner sep=0pt] (v0) at (21.6,11.75) {VPN0 = 0x003 (10)};
\node[med, minimum width=3.4cm, minimum height=0.5cm, anchor=west, inner sep=0pt] (vo) at (24.6,11.75) {offset = 0xA44 (12)};

% TLB
\node[anchor=north west, font=\scriptsize\bfseries, text=cpud] at (15.1,10.8) {TLB: 32 entries, all compared at once};
\foreach \e/\vv/\vpn/\ppn [count=\k from 0] in {0/1/0x00010/0x80020, 1/1/0x00403/0x80001, {\vdots}/{}/{}/{}, 31/0/0x7FFFF/0x00000}{
  \pgfmathsetmacro{\yy}{9.9-\k*0.5}
  \node[lab, anchor=east] at (15.6,\yy) {\e};
  \draw[pcb, fill=white] (15.7,\yy-0.2) rectangle ++(0.4,0.4);
  \draw[cpud, fill=cpuf] (16.1,\yy-0.2) rectangle ++(2.0,0.4);
  \draw[medd, fill=medf] (18.1,\yy-0.2) rectangle ++(2.0,0.4);
  \draw[pcb, fill=white] (20.1,\yy-0.2) rectangle ++(1.4,0.4);
  \node[font=\tiny\ttfamily] at (15.9,\yy) {\vv};
  \node[font=\tiny\ttfamily] at (17.1,\yy) {\vpn};
  \node[font=\tiny\ttfamily] at (19.1,\yy) {\ppn};
  \ifx\vv\empty\else\node[font=\tiny] at (20.8,\yy) {RWXU..};\fi}
\node[lab, anchor=south] at (15.9,10.1) {V}; \node[lab, anchor=south] at (17.1,10.1) {VPN (20)};
\node[lab, anchor=south] at (19.1,10.1) {PPN (22)}; \node[lab, anchor=south] at (20.8,10.1) {flags (7)};
\draw[bad, line width=1.4pt] (15.65,9.18) rectangle (21.55,9.62);
\node[note, anchor=north west, text width=10.8cm] at (15.1,5.6) {%
  50 bits $\times$ 32 entries = 1,600 bits. Small and compared in parallel $\Rightarrow$ usually \textbf{flip-flops + 32 comparators}
  (a CAM), not SRAM.\\[3pt]
  VPN 0x00403 matches entry 1 $\Rightarrow$ PPN 0x80001, in the same cycle as the cache index.};
\draw[dat] (v0.south) -- (21.6,10.75) -- (21.85,10.75) -- (21.85,9.4) -- (21.55,9.4);
\node[lab, anchor=west] at (21.9,10.2) {VPN = VPN1:VPN0};

\node[anchor=west, font=\scriptsize\bfseries] at (15.1,7.3) {physical =};
\node[med, minimum width=4.2cm, minimum height=0.5cm, anchor=west, inner sep=0pt] (pp) at (17.0,7.3) {PPN = 0x80001 (22)};
\node[med, minimum width=3.4cm, minimum height=0.5cm, anchor=west, inner sep=0pt] (po) at (21.2,7.3) {offset = 0xA44};
\node[note, anchor=north west, text width=10.8cm] at (15.1,6.75) {%
  = \textbf{0x8000\_1A44}: the address the cache above used. The offset is copied, only the page number is translated.};
\draw[dat] (19.1,9.15) -- (19.1,7.55);
\draw[dat] (vo.south) |- (22.9,8.1) -- (po.north);

% page walk
\node[anchor=north west, font=\scriptsize\bfseries, text=memd] at (26.6,10.8) {TLB miss: hardware page-table walk};
\node[pwr, minimum width=2.5cm, minimum height=0.8cm, anchor=north west, font=\tiny] (satp) at (26.6,10.3)
  {\textbf{satp} register\\root table PPN};
\node[sram, minimum width=2.6cm, minimum height=3.6cm, anchor=north west, font=\tiny, text width=2.4cm] (pt1) at (29.6,10.3)
  {\textbf{level-1 table}\\1024 PTEs $\times$ 4 B\\(in DRAM)\\[4pt]entry [VPN1 = 1]\\$\rightarrow$ next table};
\node[sram, minimum width=2.6cm, minimum height=3.6cm, anchor=north west, font=\tiny, text width=2.4cm] (pt0) at (33.0,10.3)
  {\textbf{level-0 table}\\1024 PTEs $\times$ 4 B\\(in DRAM)\\[4pt]entry [VPN0 = 3]\\$\rightarrow$ PPN 0x80001};
\draw[dat] (satp.east) -- (pt1.west |- satp.east);
\draw[dat] (pt1.east) -- (pt0.west);
\node[note, anchor=north west, text width=10.9cm] at (26.6,5.75) {%
  $\bullet$ The walker is a small FSM: 2 memory reads (each may hit in the caches), then it writes the result into a TLB entry and the access is retried.\\[2pt]
  $\bullet$ PTE V = 0 or wrong permission $\Rightarrow$ \textbf{page fault} exception, the OS fixes it.\\[2pt]
  $\bullet$ The NPU and GPU have their own MMU (an IOMMU) so they can use the same virtual addresses as the app.};

\figfoot{Sizes and latencies are typical ranges. The cache example uses a 32-bit physical address so that one address can be followed from the TLB into the cache.}
\end{tikzpicture}
\end{document}
